% AUTOMATISCH ERZEUGT aus dem SysML-Modell mit tools/generate_docs.py.
% Nicht von Hand aendern: Inhalte im Modell pflegen, Layout im Kapitel.
REQ-B-MIS-001 & Funktional & Das System muss einen Transportauftrag mit eindeutigem Start- und Zielstandort annehmen oder begruendet ablehnen. & Test: gueltiger, ungueltiger und nicht ausfuehrbarer Auftrag ueber Web-App. \\
REQ-B-UI-001 & Schnittstelle & Die erste Version muss Transportauftraege ueber eine Web-App ausloesen koennen. & Test: Bediener waehlt Start und Ziel; Roboter erhaelt Auftrag. \\
REQ-B-STA-001 & Funktional & Das System muss den Auftragsstatus mindestens als bereit, angenommen, unterwegs, Ziel erreicht, pausiert, blockiert, gestoppt oder Fehler melden. & Integrationstest mit Statusfolge je Use Case. \\
REQ-B-NAV-001 & Funktional & Das System muss einen definierten Zielstandort autonom anfahren. & Systemtest: Ziel wird ohne manuelle Fahrzeugsteuerung erreicht. \\
REQ-B-NAV-002 & Funktional & Das System muss eine Route durch mehrere Raeume, einen Flur und mindestens eine offene Tuer ausfuehren koennen. & Systemtest im Schul-Demonstrationsbereich oder Testaufbau. \\
REQ-B-NAV-003 & Funktional & Das System muss erkennen und melden, wenn Ziel oder Route nicht erreichbar sind. & Test: Route blockieren; Status blockiert oder nicht erreichbar erwarten. \\
REQ-B-LOC-001 & Funktional & Das System muss seine Pose im definierten Betriebsbereich bestimmen. & Test mit noch festzulegender Positions- und Orientierungsabweichung. \\
REQ-B-PER-001 & Funktional & Das System muss Hindernisse im relevanten Fahrbereich erfassen und fuer die lokale Navigation bereitstellen. & Test: Hindernis erscheint in lokaler Umgebungsdarstellung. \\
REQ-B-MAP-001 & Funktional & Das System muss eine freigegebene Grundkarte des Betriebsbereichs laden und fuer Navigation verwenden. & Test: Karte laden und Zielroute planen. \\
REQ-B-MAP-002 & Constraint & Das System muss temporaere Hindernisse von dauerhaften Karteninformationen unterscheiden. & Test: Hindernis blockiert Route lokal, veraendert aber die Langzeitkarte nicht automatisch. \\
REQ-B-NAV-004 & Funktional & Das System muss bei einem temporaeren Hindernis stoppen, lokal umplanen oder eine Blockade melden. & Test: Hindernis in Fahrweg platzieren; erwartete Reaktion pruefen. \\
REQ-B-SAFE-001 & Safety & Das System muss bei nicht sicher behandelbaren Hindernissen in einen sicheren Zustand uebergehen. & Test: keine Kollision; Bewegung endet kontrolliert; Status wird gemeldet. \\
REQ-B-SAFE-002 & Safety & Das System muss Personen im Fahrbereich als sicherheitsrelevante Umgebungselemente behandeln. & Safety-Testkonzept offen; zunaechst reduzierte MVP-Reaktion: stoppen und melden. \\
REQ-B-SAFE-003 & Safety & Das System muss betrieblichen Stopp, SafetyStop und EmergencyStop begrifflich und funktional trennen. Siehe DEC-013. & Review, Pause-/Fortsetzungstests und spaeterer Safety-Test; EmergencyStop bleibt Open Issue. \\
REQ-B-PERF-001 & Leistung & Die MVP-Nominalgeschwindigkeit soll ca. $0{,}5\,\mathrm{m/s}$ betragen. & Messfahrt mit noch festzulegender Toleranz. \\
REQ-B-PWR-001 & Nichtfunktional & Die erste Version soll ein definiertes Unterrichtsprofil von ca. 2 h abdecken. & Dauertest mit noch zu definierendem Fahr-/Warte-/Demonstrationsprofil. \\
REQ-B-TRN-001 & Schnittstelle & Der Roboter muss einen manuell aufgesetzten Werkstuecktraeger als Payload mitfuehren koennen. & Mechanischer Test offen, abhaengig von OI-MECH-001 und OI-MECH-002. \\
REQ-B-ENE-001 & Funktional & Das System muss eine Rueckkehr zur Ladestation als Ziel oder Betriebsbefehl behandeln koennen. & MVP-Test: Fahrziel Ladestation erreichen; Docking/Laden offen. \\
